\documentclass[10pt,a4paper]{amsart}
\usepackage[margin=19mm]{geometry}
\usepackage{amsmath,amssymb,mathtools}
\usepackage[scr=rsfs,cal=euler]{mathalfa}
\usepackage{xfrac}
\usepackage[unicode,hidelinks,bookmarks=false]{hyperref}
\newcommand{\ie}{i.e.\ }
\newcommand{\cf}{cf.\ }
\newcommand{\resp}{respectively\ }
\newcommand{\Subsection}{\subsection}
\providecommand{\nobreakdash}{\nobreak\hbox{-}}
\setlength{\emergencystretch}{2em}
\multlinegap0pt
\usepackage{bm}
\usepackage{amssymb}
\usepackage{xcolor}
\usepackage{algorithm}\usepackage{algpseudocode}
\usepackage{tikz}
\usepackage{mathtools}
\usepackage{cancel}
\newtheorem{proposition}{Proposition}
\newcommand{\new}[1]{#1}
\title{A relax-fix-and-exclude algorithm for an MINLP problem with multilinear interpolations}
\author{Bruno Machado Pacheco, Pedro Marcolin Antunes, Eduardo Camponogara, Laio Oriel Seman, Vin\'icius Ramos Rosa, Bruno Ferreira Vieira, Cesar Longhi}
\date{Source: arXiv:2502.21249v2 (CC BY 4.0)}
\begin{document}
\maketitle

\noindent\emph{Source and licence.} A relax-fix-and-exclude algorithm for an MINLP problem with multilinear interpolations. Version arXiv:2502.21249v2 (CC BY 4.0), distributed by arXiv under the Creative Commons Attribution 4.0 International licence, http://creativecommons.org/licenses/by/4.0/, as recorded in the arXiv licence metadata for this exact version and shown on the abstract page https://arxiv.org/abs/2502.21249v2. Full licence notice: \texttt{licenses/arxiv-2502.21249-CC-BY-4.0.txt}.\par\smallskip

\noindent\textit{Editorial scope.} Counted unit: the article's proposition prop (source0.tex lines 832--834). The article's complete original proof of it is delivered with it (source0.tex lines 835--859, one \texttt{proof} environment of 25 lines). No mathematics is written, repaired or re-derived: the statement and the proof are the article's own lines, in the article's own notation, and no retained line is substituted or re-flowed.  Retained uncounted as the source-local context the proof needs: lines 582--584; lines 594--609; lines 800--802; lines 805--820.  Nothing was omitted from any retained span, so the package carries no omission record.  No retained line contains a citation, so the wrapper carries no bibliography and no bibliography entry.  No retained line contains a figure environment, an image-inclusion command or drawing code, so no image is delivered and the package declares no asset dependency.  Editorial formatting only, in addition: an AMS \texttt{amsart} preamble that restates the article's own declarations and macros used by the retained lines, namely the environments \texttt{proposition} on separate counters, together with the standard packages those lines need.  The retained environments print in this document as Proposition 1 in order of appearance, whereas the article prints the same items with the numbering of its own full document.  The complete native source of the article as distributed in the arXiv e-print for this exact version is bundled unchanged as \texttt{sources/174048/source0.tex}.
\begin{equation}\label{eq:interpolation-bilinears}
   \lambda_{\bm{k}} = \prod_{j=1}^{n} \xi^j_{k_j},\quad \forall\bm{k} \in \mathcal{I}
.\end{equation}

\begin{equation}\label{eq:minlp}\tag{$\mathcal{P}$}
\begin{aligned}
    C = \min_{\bm{\xi},\bm{y}} \quad & \widetilde{f}_0 + \bm{c}^{T} \bm{y}  \\
    \textrm{s.t.} \quad & \widetilde{f}_i + \bm{a}_i^{T} \bm{y} \le b_i, & i=1,\ldots,m, \\
                      %%%
			& \widetilde{f}_i = \sum_{\bm{k}\in\mathcal{I}} \lambda_{\bm{k}}  f_i(\hat{\bm{x}}^{\bm{k}}), & i=0,\ldots,m, \\
                   %%%
   &\lambda_{\bm{k}} = \prod_{j=1}^{n} \xi^j_{k_j}, &
     \forall \bm{k}\in\mathcal{I}, \\
       %%%%
   &x_{j} = \sum_{k\in\mathcal{I}_j} \xi^j_{k} \hat{x}_{j}^{k}, &j=1,\ldots,n, \\
   &1 = \sum_{k\in\mathcal{I}_j} \xi^j_{k}, &j=1,\ldots,n, \\
      & \bm{\xi^j} \in \text{SOS2}(\mathcal{I}_j), \bm{\xi^j}\geq \mathbf{0} & j=1,\ldots,n, \\
      & \bm{x}\in X, \,\bm{y}\in \left\{ 0,1 \right\}^{p}
,\end{aligned}
\end{equation}

\begin{equation}\label{eq:relaxed-interpolation-bilinears}
     \xi^j_{\hat{k}_j} = \sum_{\substack{\bm{k} \in \mathcal{I}: \\ k_j = \hat{k}_j}} \lambda_{\bm{k}},\quad \forall \hat{k}_j\in \mathcal{I}_j,\, j=1,\ldots,n .
\end{equation}

\begin{equation}\label{eq:milp}\tag{$\overline{\mathcal{P}}$ }
\begin{aligned}
    \overline{C} = \min_{\bm{\xi},\bm{\lambda},\bm{y}} \quad & \widetilde{f}_0 + \bm{c}^{T} \bm{y}  \\
    \textrm{s.t.} \quad & \widetilde{f}_i + \bm{a}_i^{T} \bm{y} \le b_i, & i=1,\ldots,m \\
		  %%%%
			& \widetilde{f}_i = \sum_{\bm{k}\in\mathcal{I}} \lambda_{\bm{k}}  f_i(\hat{\bm{x}}^{\bm{k}}), & i=0,\ldots,m, \\
                   %%%
        &\xi^j_{\hat{k}_j} = \sum_{\substack{\bm{k} \in \mathcal{I}: \\ k_j = \hat{k}_j}} \lambda_{\bm{k}},& \forall \hat{k}_j\in \mathcal{I}_j,\, j=1,\ldots,n, \\
       %%%%
   &x_{j} = \sum_{k\in\mathcal{I}_j} \xi^j_{k} \hat{x}_{j}^{k}, &j=1,\ldots,n, \\
   &1 = \sum_{k\in\mathcal{I}_j} \xi^j_{k}, &j=1,\ldots,n, \\
      & \bm{\xi^j} \in \text{SOS2}(\mathcal{I}_j), & j=1,\ldots,n, \\
      & \bm{\lambda} \in [0,1]^{|\mathcal{I}|}, \\
      & \bm{x}\in X, \,\bm{y}\in \left\{ 0,1 \right\}^{p} 
.\end{aligned}
\end{equation}

\begin{proposition}\label{prop}
    Problem \ref{eq:milp} is a relaxation of \ref{eq:minlp}.
\end{proposition}

\begin{proof}
    Because the only change from \ref{eq:minlp} to \ref{eq:milp} is the removal of constraint \eqref{eq:interpolation-bilinears} and the addition of constraint \eqref{eq:relaxed-interpolation-bilinears}, it suffices to show that any solution of \ref{eq:minlp} satisfies \eqref{eq:relaxed-interpolation-bilinears}.
    Thus, let $( \widetilde{\bm{x}}, \widetilde{\bm{y}}, \widetilde{\bm{\xi}},\widetilde{\bm{\lambda}} ) $ be a feasible solution of \ref{eq:minlp}.
    \new{%
    Consider first the case $j=1$ and take any $\hat{k}_1\in\mathcal{I}_1$.
    Because $\widetilde{\lambda}_{\bm{k}}$ respects \eqref{eq:interpolation-bilinears} for any $\bm{k}\in\mathcal{I}$, we have
    \begin{equation}
    \begin{aligned}
        \sum_{\substack{\bm{k} \in \mathcal{I}: \\ k_1 = \hat{k}_1}} \widetilde{\lambda}_{\bm{k}}
            &= \sum_{\substack{\bm{k} \in \mathcal{I}: \\ k_1 = \hat{k}_1}} \prod_{j=1}^{n} \widetilde{\xi}^j_{k_j} \\
                  %%%%%%
            &= \widetilde{\xi}^1_{\hat{k}_1} \sum_{k_2 \in \mathcal{I}_2} \widetilde{\xi}^2_{k_2} \sum_{k_3 \in \mathcal{I}_3} \widetilde{\xi}^3_{k_3} \cdots \sum_{k_n \in \mathcal{I}_n}  \widetilde{\xi}^n_{k_n} \\
                     %%%%%
            &= \widetilde{\xi}^1_{\hat{k}_1},
    \end{aligned}
    \end{equation}
    where the last equality comes from $\widetilde{\bm{\xi}}$ satisfying
    \begin{equation}
	\sum_{k_j\in \mathcal{I}_j} \widetilde{\xi}^j_{k_j} = 1,\,j=1,\ldots,n
    \end{equation}
    as it is feasible for \ref{eq:minlp}.
    Therefore, the first constraint of \eqref{eq:relaxed-interpolation-bilinears} is satisfied, given that the choice for $\hat{k}_1$ was arbitrary.
    \emph{Mutatis mutandis}, the same can be shown for $j=2,\ldots,n$, implying that all constraints in \eqref{eq:relaxed-interpolation-bilinears} are satisfied by a feasible solution of \ref{eq:minlp}.
    }
\end{proof}

\end{document}
